% Appendix D. Fragment, \input by spacenav.tex inside \appendix.
% Commands transcribed from the scripts' own usage lines; versions from
% requirements-gpu.txt and the conda environment the runs were made in.

\section{Reproduction}
\label{app:repro}

Everything reported in this paper is produced by scripts in the repository,
from a checkout with the \code{spacenav} package installed and the GPU build
of JAX, in the order below; \code{data/} is not under version control, and the
level pools are regenerated from their seed blocks (seeds from 0 for the
training pool, from $10^7$ for the seen-family validation pool, from
$2 \times 10^7$ for the holdout pool, from $3 \times 10^7$, $4 \times 10^7$
and $5 \times 10^7$ for the three long-horizon pools: \code{long}, the
\code{long\_val} pool the 600\,s suites are drawn from, and the 900-level pool
validated to 9{,}000\,s) so that the levels themselves
are reproducible without being distributed.

\begin{small}
\begin{verbatim}
# level pools: train, val_seen, val_holdout (data/pools/*.npz), then the
# long-horizon pool used at a 600 s clock
python scripts/build_pools.py --train 2000 --val 300 --holdout 400
SPACENAV_EPISODE_S=600 python scripts/build_pools.py --only long --n 200

# one training run at the 150 s clock; the frozen suites
# (data/eval/suites.pkl) are built on first use
SPACENAV_INSTRUMENTS=eta,stretch \
  python scripts/train.py --run abl_info_s0 --updates 3000 --seed 0 --pool train
SPACENAV_INSTRUMENTS=0 \
  python scripts/train.py --run ppo_v5_noinstr --updates 3000 --seed 0
SPACENAV_INSTRUMENTS=1 \
  python scripts/train.py --run ppo_v4 --updates 3000 --seed 0
SPACENAV_INSTRUMENTS=budget,demand \
  python scripts/train.py --run abl_redundant_s0 --updates 3000 --seed 0

# the 600 s clock: suites from the long_val pool (built here on first use),
# then training on the long pool against them
SPACENAV_EPISODE_S=600 python scripts/build_long_suites.py --clock 600
SPACENAV_EPISODE_S=600 SPACENAV_SUITES=data/eval/suites_long.pkl \
  python scripts/train.py --run long600_s0 --updates 3000 --seed 0 --pool long

# the clock sweep (the reference optimizer at 150, 300, 600 and 900 s)
bash scripts/clock_sweep.sh data/curriculum/clock

# the ghost fleets: every checkpoint of one run on one instance, rolled out
# with the instrument setting the run was trained with
SPACENAV_INSTRUMENTS=eta,stretch python scripts/ghost_fleet.py \
  --run abl_info_s0 --suite medium --task 1 --rollouts 16 \
  --out data/ghosts/medium_1.json

# the zero-thrust navigator to 900 s, the 9,000 s pool, and the same to 9,000 s
python scripts/noop_clock.py measure --clock 900
python scripts/noop_clock.py build10x --per-family 100
python scripts/noop_clock.py measure10x --clock 9000
python scripts/noop_clock.py figure

# per-instance success against eta, A, D and family
python scripts/correlations.py

# the training-budget run
SPACENAV_INSTRUMENTS=eta,stretch \
  python scripts/train.py --run budget10k_s0 --updates 10000 --seed 0

# figures
python scripts/paper_figures.py --which ghosts,ablation,clock,curves,eta
python scripts/fig_task.py --suite hard --task 63 --run ppo_v5_noinstr
python scripts/fig_phase.py compute --pool data/pools/long.npz
python scripts/fig_phase.py render
\end{verbatim}
\end{small}

Two environment variables select the process-wide settings that the body of
the paper varies. \code{SPACENAV\_EPISODE\_S} sets the episode clock in
seconds (default \code{150.0}) and is read once, at import, by
\code{spacenav/constants.py}, from which the tick count, the time-remaining
entries of the observation and the number of snapshots a pool stores are all
derived; a clock longer than the window a pool was validated over (300\,s for
\code{train} and the two validation pools, 900\,s for the long-horizon pools)
runs late-starting episodes past where the level was checked, which is why the
long clock has its own pools and its own suites, selected with
\code{SPACENAV\_SUITES}, and why \code{train.py} refuses to build a
non-default suite file on its own. \code{SPACENAV\_INSTRUMENTS} selects which
of the four instrument channels are live (Appendix~\ref{app:obs}), with the
training default \code{eta,stretch}, and since every setting produces the same
observation shape a checkpoint must be evaluated or replayed under the setting
recorded in its run's \code{config.json}, as nothing will complain if it is
not. Every evaluation reported here is the greedy policy (the most probable
action at each tick) on the frozen suites, written by \code{train.py} to
\code{eval.jsonl} at every fiftieth update, with the quoted success rate being
the mean of a run's last three checkpoints; the ghost-fleet replays are the
one place actions are sampled.

The simulation, environment, optimizer and training loop are written in
Python 3.12 against JAX 0.9.2 (the \code{jax[cuda12]} build), with the policy
network in Flax 0.12.6 and the optimizer state in Optax 0.2.8; array handling
and the statistics in the figure scripts use NumPy 2.4.3 and SciPy 1.17.1, and
the figures are drawn with matplotlib 3.11.2, all as pinned in
\code{requirements-gpu.txt} and installed in the conda environment the runs
were made in. The browser client in \code{web/} is a direct port of the
physics in JavaScript, and a parity test (\code{tests/test\_js\_parity.py})
holds the two integrators to the same trajectories, which is what makes the
replays on the companion site the flights that were scored rather than
re-renderings of them.

The twenty headline runs were made on two cards. Seeds 0 and 1 of the
no-channel and all-four arms ran on an NVIDIA GeForce GTX 1080 Ti (11\,GB,
Pascal), at about 17,000 environment steps per second averaged over a whole
run including its checkpoint evaluations (the training loop alone measures
31,900 steps/s at this network width), or 6.4\,h for 3,000 updates; the
remaining sixteen runs, and all of the long-clock runs, ran on an NVIDIA
GeForce RTX 3090 (24\,GB) at about 111,000 steps per second for a single
process, or 67\,min per run, a card that one such process already saturates,
so that two concurrent runs together deliver three quarters of the throughput
of one. Within the no-channel arm, whose seeds are split across the two cards, no
difference attributable to the card was found. The cross-entropy optimizer, the chaos probe and the phase portraits run
on either card, with the single-burn outcome maps of Figure~\ref{fig:phase}
being the one computation (368,640 full-episode rollouts per level) that is
impractical on the smaller one.

The reasoning behind every design choice, rejected approach and measurement
that changed the plan is recorded in \code{docs/decisions/}, as numbered,
append-only decision records; a reversed decision is superseded by a new
record rather than edited, so that the wrong turns remain legible, and each
record cites the script, sample size or paper that forced it. The records are
the trail from which the arguments of Sections~\ref{sec:eta}
through~\ref{sec:results} were assembled (for instance, the equal-pull surface
replacing the resonance-overlap scale, the solvability yardstick being judged
on sample size, the instrument channels being split along the redundancy
line, and the clock becoming a per-process setting), and they are the place
to look when a number in this paper and a number in the repository disagree,
since the record will say which measurement superseded which.
